\documentclass[10pt]{article}
% Math Packages
\usepackage{amsmath, mathtools}
\usepackage{amssymb}
\usepackage{amsthm}
\usepackage{amsfonts}
\usepackage{bbm}
\usepackage{breqn}
\usepackage[margin=1in]{geometry}
\usepackage{graphicx}
\usepackage{tikz}
\usetikzlibrary{arrows.meta}
\usetikzlibrary{calc}
\usepackage{forest}
\usepackage{tikz-qtree}
\graphicspath{ {./images/} }
\usepackage{hyperref}
\usepackage[capitalize]{cleveref}
\usepackage[shortlabels]{enumitem}
\usetikzlibrary{arrows,matrix,positioning}
\usepackage{multicol}

% for the pipe symbol
\usepackage[T1]{fontenc}

% Citing theorems by name. (source: https://tex.stackexchange.com/questions/109843/cleveref-and-named-theorems)
\makeatletter
\newcommand{\ncref}[1]{\cref{#1}\mynameref{#1}{\csname r@#1\endcsname}}

\def\mynameref#1#2{%
  \begingroup
  \edef\@mytxt{#2}%
  \edef\@mytst{\expandafter\@thirdoffive\@mytxt}%
  \ifx\@mytst\empty\else
  \space(\nameref{#1})\fi
  \endgroup
}
\makeatother

% Colorful Notes
\usepackage{color} \definecolor{Red}{rgb}{1,0,0} \definecolor{Blue}{rgb}{0,0,1}
\definecolor{Purple}{rgb}{.5,0,.5} \def\red{\color{Red}} \def\blue{\color{Blue}}
\def\gray{\color{gray}} \def\purple{\color{Purple}}
\newcommand{\rnote}[1]{{\red [#1]}} % \rnote{foo} gives '[foo]' in red
\newcommand{\pnote}[1]{{\purple [#1]}} % \pnote{foo} gives '[foo]' in purple
\newcommand{\bnote}[1]{{\blue #1}} % \bnote{foo} gives 'foo' in blue
\newcommand{\gnote}[1]{{\gray #1}} % \gnote{foo} gives 'foo' in gray
\newcommand{\Max}[1]{{\purple [#1]}} % \bnote{foo} then 'foo' is blue


% Claim numbering (the counter restarts after each proof environment)
\newcounter{claimcount}
\setcounter{claimcount}{0}
\newenvironment{claim}{\refstepcounter{claimcount}\par\addvspace{\medskipamount}\noindent\textbf{Claim \arabic{claimcount}:}}{}
\usepackage{etoolbox}
\AtBeginEnvironment{proof}{\setcounter{claimcount}{0}}
\newenvironment{claimproof}{\par\addvspace{\medskipamount}\noindent\textit{Proof of Claim  \arabic{claimcount}.}}{\hfill\ensuremath{\qedsymbol} \tiny{Claim}

  \medskip}
% Add claim support to cleverref
\crefname{claimcount}{Claim}{Claims}


% Math Environments
\newtheorem{theorem}{Theorem}
\newtheorem{assumption}[theorem]{Assumption}
\newtheorem{lemma}[theorem]{Lemma}
\newtheorem{proposition}[theorem]{Proposition}
\newtheorem{corollary}[theorem]{Corollary}
\newtheorem{question}[theorem]{Question}
\theoremstyle{definition}
\newtheorem{definition}[theorem]{Definition}
\newtheorem{remark}[theorem]{Remark}
\newtheorem{example}[theorem]{Example}
\newtheorem{notation}[theorem]{Notation}
\newtheorem{problem}[theorem]{Problem}

% Matrices and Column Vectors. 
\usepackage{stackengine}
\setstackgap{L}{1.0\normalbaselineskip}
\usepackage{tabstackengine}
\setstackEOL{;}% row separator
\setstackTAB{,}% column separator
\setstacktabbedgap{1ex}% inter-column gap
\setstackgap{L}{1.5\normalbaselineskip}% inter-row baselineskip
\let\nmatrix\bracketMatrixstack  %Usage: \nmatrix{1,2,3\4,5,6}
\newcommand\cv[1]{\setstackEOL{,}\bracketMatrixstack{#1}} %usage: \cv{1,2,3}

% Custom Math Coqmmands
\newcommand{\vt}{\vskip 5mm} % vertical space
\newcommand{\fl}{\noindent\textbf} % first line
\newcommand{\Fl}{\vt\noindent\textbf} % first line with space above
\newcommand{\norm}[1]{\left\lVert#1\right\rVert} % norm
\newcommand{\pnorm}[1]{\left\lVert#1\right\rVert_p} % p-norm
\newcommand{\qnorm}[1]{\left\lVert#1\right\rVert_q} % q-norm
\newcommand{\1}[1]{\textbf{1}_{\left[#1\right]}} % indicator function 
\def\limn{\lim_{n\to\infty}} % shortcut for lim as n-> infinity
\def\sumn{\sum_{n=1}^{\infty}} % shortcut for sum from n=1 to infinity
\def\sumkn{\sum_{k=1}^{n}} % shortcut for sum from k=1 to n
\def\sumin{\sum_{i=1}^{n}} % shortcut for sum from i=1 to n
\def\SAs{\sigma\text{-algebras}} % shortcut for $\sigma$-algebras
\def\SA{\sigma\text{-algebra}} % shortcut for $\sigma$-algebra
\def\Ft{\mathcal{F}_t} % time-indexed sigma-algebra (t)
\def\Fs{\mathcal{F}_s} % time-indexed sigma-algebra (s)
\def\F{\mathcal{F}} % sigma-algebra
\def\G{\mathcal{G}} % sigma-algebra
\def\R{\mathbb{R}} % Real numbers
\def\N{\mathbb{N}} % Natural numbers
\def\Z{\mathbb{Z}} % Integers
\def\E{\mathbb{E}} % Expectation
\def\P{\mathbb{P}} % Probability
\def\Q{\mathbb{Q}} % Q probability
\def\dist{\text{dist}} %Text 'dist' for things like 'dist(x,y)'
\newcommand{\indep}{\perp \!\!\! \perp}  %independence symbol
\def\Var{\mathrm{Var}} % Variance
\def\tr{\mathrm{tr}} % trace

% Brackets and Parentheses
\def\[{\left [}
    \def\]{\right ]}
% \def\({\left (}
%   \def\){\right )}



\usepackage{color}
\definecolor{Red}{rgb}{1,0,0}
\definecolor{Blue}{rgb}{0,0,1}
\definecolor{Purple}{rgb}{.5,0,.5}
\def\red{\color{Red}}
\def\blue{\color{Blue}}
\def\gray{\color{gray}}
\def\purple{\color{Purple}}
\definecolor{RoyalBlue}{cmyk}{1, 0.50, 0, 0}
\newcommand{\dempfcolor}[1]{{\color{RoyalBlue}#1}} 
\newcommand{\demph}[1]{\dempfcolor{{\sl #1}}}

% comment exactly one of the following line to show / hide the solutions
% \newcommand{\solution}[1]{{\purple #1}} % uncomment to show the solutions
\newcommand{\solution}[1]{} % uncomment to hide the solutions



\title{Lecture Notes for Math 372: \\Elementary Probability and Statistics}
\date{Last updated: \today}
% \author{mh}

\begin{document}
\begin{center}
  \section*{Math 372: Homework 12}
  \textit{Due Wednesday, May 6}
\end{center}

\textit{Instructions: Do not write your solutions on this sheet---you don't have enough space to show
your work. So use separate sheets of paper. For all problems you need to show your work.}




\begin{problem}%[best-fit line]
  Consider the following three data points, of the form $(x,y)$:
  \begin{equation*}
    (-1,4), \quad (0,5), \quad \text{and} \quad (1,9)
  \end{equation*}
  \begin{enumerate}[(a)]
    \item Plot the data points on the $xy$-plane
    \item Suppose we have a line $y=C+Dx$. Let
    $\vec{e}= (e_{1},e_{2},e_{3})$ be the vector
    whose entries are the vertical distances from the data points to the line.
    Write out $\vec{e}$.
    \item Find the values of $C$ and $D$ which minimize the length of
    $\norm{\vec{e}}^{2}=e_{1}^{2}+e_{2}^{2}+e_{3}^{2}$.
    \item Sketch the line $y=C+Dx$ using your values of $C$ and $D$.
  \end{enumerate}
\end{problem}


\begin{problem}[Salk Polio Field Trial (1954)]
  \label{ex:salk-polio-field-trial-1954}
  The results of a 1954 polio vaccine trial in $n=401,974$ elementary school
  students gave the following results:
  \begin{equation}\label{eq:table-of-observed-values}
    \begin{tabular}{l|cc}
      & \textbf{No Polio} & \textbf{Polio} \\
      \hline
      \textbf{Vaccinated} & 200,712          & 33            \\
      \textbf{Unvaccinated}  & 201,114           & 115           
    \end{tabular} 
  \end{equation}
  We wish to know if there is a relationship between vaccination status and
  the incidience of paralytic polio. The null hypothesis is:
  \begin{equation*}
    H_{0}= \text{vaccination and polio are independent}
  \end{equation*}
  and the alternative hypothesis is
  \begin{equation*}
    H_{a} = \text{ there is a relationship between vaccination and polio.}
  \end{equation*}
  Conduct a hypothesis test using the chi-squared statistic to test these
  hypotheses.
\end{problem}

% \begin{problem}
%   The time in hours $T$ between streptococcal infection and display of a sore throat is a
%   random variable whose probabililty density function is approximately
%   \begin{equation*}
%     f(t) 
%     = \left\{ 
%       \begin{array}{l@{\quad:\quad}l} \frac{1}{15676} t^{2} e^{-0.05 t},  &0 \leq t \leq 150\\ 0& \text{else} \end{array}\right.
%   \end{equation*}
%   \begin{enumerate}[(a)]
%     \item What is the probability that an infected patient will display
%     symptoms within the first 48 hours?
%     \item What is the probability that an infected patient will not display
%     symptoms until after 36 hours?
%   \end{enumerate}
% \end{problem}


\begin{problem}
  Let $S_{200}$ be the number of heads that turn up in 200 tosses of a fair
  coin. Estimate
  \begin{equation*}
    \P\left[S_{200}=100 \right], \P\left[S_{200}=90 \right], \text{ and } \P\left[S_{200}=80 \right].
  \end{equation*}
\end{problem}



\begin{problem}
  Find the probability that among $10,000$ random digits the digit 3 appears
  not more than 931 times.
\end{problem}


\begin{problem}
  A random walker starts at 0 on the x-axis and at each time unit moves 1 step
  to the right or 1 step to the left with probability 1/2. Estimate the
  probability that, after 100 steps, the walker is more than 10 steps from the
  starting position.
\end{problem}



\begin{problem}
  An urn contains one red ball and one blue ball. A ball is drawn at random.
  If it is red, the game is over. If it is blue, it is replaced in the urn
  along with an extra blue ball. This process continues until the red ball is
  drawn, however long it takes. Let $Y$ be the number of draws required.
  \begin{enumerate}[(a)]
    \item Argue that $\P\left[Y=k \right] = \frac{1}{k(k+1)}$.
    \item Show that this defines a valid pmf. (Hint: first show that
    $\frac{1}{k(k+1)}= \frac{1}{k}-\frac{1}{k+1}$, and then consider
    $\sum_{k=1}^{n}\frac{1}{k(k+1)}$).
    \item Show that $\E\left[Y \right]$ does not exist.
  \end{enumerate}1
\end{problem}

\newpage

\begin{problem}%[best-fit-line]
  Consider the following four datapoints:
  \begin{equation*}
    (-1,2), \quad (0,0), \quad (1,-3), \quad\text{and} \quad (2,-5)
  \end{equation*}
  \begin{enumerate}[(a)]
    \item Find the best straight-line fit to the data. 
    \item Plot points and the line.
  \end{enumerate}
\end{problem}


\begin{problem}[Normal Distribution]
  If $Z\sim \mathcal{N}(0,1)$, then
  \begin{equation}\label{eq:1}
    \P\left[|Z| \leq u \right] = 2 \P\left[Z \leq u\right]-1
  \end{equation}
  for any $u \geq 0$. Draw a picture to illustrate the geometric meaning of \cref{eq:1}.
\end{problem}


\end{document}